
\chapter{Introduzione}\label{chap:introduction}

% ─────────────────────────────────────────────────────────────────────────────
\section[Digitalizzazione, minacce e centralità della Digital Forensics]{Digitalizzazione, minacce e centralità della \glsentrylong{df}}
% ─────────────────────────────────────────────────────────────────────────────
La progressiva digitalizzazione della società contemporanea ha trasformato in profondità le modalità con cui le informazioni vengono prodotte, archiviate, trasmesse e analizzate. Infrastrutture cloud, servizi distribuiti, dispositivi mobili interconnessi, piattaforme sociali e sistemi informativi eterogenei generano oggi una quantità crescente di dati digitali, che assumono un ruolo sempre più rilevante sia nella vita quotidiana sia nelle attività professionali, istituzionali e giudiziarie. Tale evoluzione ha ampliato in modo significativo anche la superficie d'attacco dei sistemi informatici, aumentando le opportunità di sfruttamento da parte di soggetti ostili e moltiplicando i possibili punti di compromissione delle evidenze digitali \cite{garfinkel2010digital,casey2011digital}.

Parallelamente alle opportunità offerte dall'innovazione tecnologica, si sono affermate minacce sempre più sofisticate, caratterizzate da elevata eterogeneità tecnica e da una crescente capacità di occultamento. Tra queste rientrano l'esfiltrazione e la manipolazione dei dati, l'alterazione di file multimediali, la cancellazione selettiva delle tracce, l'impiego di tecniche di anonimizzazione e, più in generale, tutte quelle pratiche che mirano a ostacolare la ricostruzione attendibile degli eventi. In un simile scenario, la \gls{df} assume un ruolo centrale quale disciplina deputata all'identificazione, acquisizione, preservazione, analisi e presentazione dell'evidenza digitale secondo criteri di rigore metodologico, integrità e verificabilità \cite{palmer2001road,reith2002examining,nist2006guide,iso27037}.

Secondo la definizione comunemente richiamata in letteratura, la \gls{df} costituisce l'applicazione di metodi scientifici e procedure tecniche finalizzate a rendere i dati digitali utilizzabili come elementi di prova in un contesto investigativo o giudiziario \cite{palmer2001road}. Tale impostazione richiede che ogni fase del processo sia condotta in modo tracciabile, ripetibile e difendibile, affinché il risultato dell'analisi non sia soltanto tecnicamente corretto, ma anche sostenibile sotto il profilo probatorio. In termini operativi, il processo forense viene comunemente ricondotto a una sequenza di fasi che comprende identificazione, acquisizione, conservazione, analisi e presentazione dell'evidenza, tutte accomunate dall'esigenza di garantire integrità, trasparenza e verificabilità \cite{reith2002examining,nist2006guide}. Questa esigenza diventa particolarmente rilevante nei casi in cui le evidenze risultino distribuite su molteplici dispositivi, conservate in ambienti virtualizzati, presenti in forma volatile oppure esposte a manipolazioni intenzionali \cite{nist2006guide,iso27037}.

La crescente incidenza dei reati informatici e, più in generale, il ruolo ormai pervasivo delle tecnologie digitali nella commissione e nella documentazione dei reati hanno consolidato la funzione della \gls{df} come snodo essenziale tra competenza tecnica e accertamento dei fatti. In questo quadro, la qualità delle procedure di analisi e l'affidabilità degli strumenti impiegati incidono direttamente sulla possibilità di ricostruire correttamente gli eventi, individuare i responsabili e preservare il valore dell'evidenza nel procedimento giudiziario \cite{costantini2019digital,clough2015principles,europol2019iocta}.

L'integrazione di sistemi basati su \gls{ai} nei workflow di \gls{df} risponde alla necessità di esaminare volumi crescenti di dati digitali e di selezionare con tempestività contenuti potenzialmente rilevanti per l'indagine. Quando tale supporto riguarda la classificazione automatizzata di immagini, tuttavia, l'affidabilità dell'analisi può essere compromessa da input degradati o intenzionalmente manipolati, capaci di ridurre la probabilità che elementi di interesse investigativo vengano correttamente segnalati all'analista. Muovendo da questa criticità, la presente tesi valuta la robustezza operativa di modelli e strumenti AI-based impiegati per la classificazione di immagini raffiguranti armi, confrontandone il comportamento su dati non manipolati, contenuti fuori distribuzione e immagini sottoposte a perturbazioni adversariali o trasformazioni anti-forensi. L'obiettivo è contribuire alla valutazione dell'automazione come supporto controllato all'analisi forense, mantenendo centrali la supervisione umana, la tracciabilità delle procedure e l'interpretazione prudente degli output automatici \cite{costantini2019digital,sanna2024preliminary,sanna2025pixelperfect}.


% ─────────────────────────────────────────────────────────────────────────────
\section[Automazione dell'analisi forense e ruolo dell'AI]{Automazione dell'analisi forense e ruolo dell'\glsentryshort{ai}}
% ─────────────────────────────────────────────────────────────────────────────
L'espansione quantitativa e qualitativa dei dati oggetto di indagine ha reso progressivamente meno sostenibile un approccio esclusivamente manuale all'analisi forense. Le metodologie tradizionali, pur mantenendo un ruolo fondamentale in termini di controllo tecnico e validazione dell'evidenza, risultano spesso difficilmente scalabili quando l'investigatore deve esaminare grandi volumi di immagini, video, log, conversazioni, file di sistema e artefatti eterogenei provenienti da fonti differenti. In tali condizioni, l'automazione rappresenta una necessità operativa prima ancora che una scelta di efficienza \cite{garfinkel2010digital,costantini2019digital}.

In questo contesto, l'integrazione di tecniche di \gls{ai} negli strumenti di \gls{df} ha introdotto nuove possibilità di supporto all'analisi investigativa. In particolare, i metodi di \gls{ml} e \gls{dl} consentono di automatizzare attività quali la classificazione semantica di immagini e video, il riconoscimento di pattern ricorrenti, il raggruppamento di evidenze simili, l'analisi di log e l'individuazione di contenuti potenzialmente rilevanti ai fini dell'indagine. La funzione di triage assume ulteriore rilievo nei casi in cui l'analista debba esaminare ripetutamente contenuti sensibili: studi condotti su professionisti impegnati nell'analisi di materiale di abuso sessuale su minori evidenziano il possibile impatto psicologico dell'esposizione protratta a tali contenuti e l'interesse operativo per tecnologie di filtering e safer presentation volte a ridurre l'esposizione non necessaria dell'operatore \cite{sanchez2019practitioner}. Pur non coincidendo con il task di classificazione \texttt{weapon}/\texttt{non\_weapon} valutato nella presente tesi, tale evidenza mostra che l'affidabilità dell'automazione può incidere non soltanto sull'efficienza del workflow, ma anche sulle condizioni operative in cui l'analista svolge la revisione dei contenuti digitali \cite{costantini2019digital,sanna2024preliminary}.


L'applicazione dell'\gls{ai} alla \gls{df} non si limita tuttavia a una mera accelerazione dei processi esistenti. Essa introduce un cambiamento qualitativo nel modo in cui l'informazione investigativa viene filtrata, ordinata e priorizzata. Sistemi basati su \gls{ai} integrati in strumenti forensi commerciali possono suggerire categorie semantiche, supportare il clustering automatico di contenuti multimediali, evidenziare anomalie e proporre correlazioni che orientano l'attenzione dell'operatore verso specifici elementi dell'evidenza. Tale tendenza è osservabile anche nei principali strumenti professionali impiegati nel settore, nei quali moduli basati su \gls{ai} vengono progressivamente integrati per supportare attività di classificazione, ordinamento e selezione preliminare delle evidenze \cite{sanna2025cybercrime,sanna2025pixelperfect}.

Tuttavia, il vantaggio offerto dall'automazione non può essere valutato esclusivamente in termini di efficienza. In ambito forense, infatti, ogni riduzione del carico umano ottenuta mediante modelli automatici deve essere bilanciata con esigenze di affidabilità, trasparenza e controllo. Un sistema che classifica automaticamente un contenuto come rilevante o irrilevante non si limita a velocizzare il lavoro dell'analista, ma contribuisce a determinare quali evidenze riceveranno attenzione prioritaria e quali, al contrario, rischieranno di essere trascurate. Ne deriva che l'introduzione dell'\gls{ai} nei workflow forensi richiede una riflessione critica non solo sulle prestazioni dei modelli, ma anche sulla solidità delle loro decisioni in scenari operativi complessi e sulla possibilità di renderne comprensibili i limiti attraverso approcci di \gls{xai} \cite{Adadi2018,Doshi-Velez2017,EUAI2020}.

% ─────────────────────────────────────────────────────────────────────────────
\section[Robustezza dei sistemi AI-based e minacce adversariali e anti-forensi]{Robustezza dei sistemi \glsentryshort{ai}-based e minacce adversariali e anti-forensi}
% ─────────────────────────────────────────────────────────────────────────────
L'impiego di sistemi basati su \gls{ai} in contesti ad alta criticità ha evidenziato una questione fondamentale: prestazioni elevate in condizioni standard non implicano necessariamente affidabilità in presenza di input alterati, degradati o ostili. Nel caso della \gls{df}, tale aspetto assume particolare rilievo, poiché i sistemi di classificazione automatizzata possono essere esposti a contenuti intenzionalmente manipolati, a trasformazioni introdotte lungo il ciclo di vita dell'immagine oppure a input che si discostano in modo significativo dalla distribuzione dei dati utilizzati in fase di sviluppo o addestramento.

Il problema dell'alterazione intenzionale dell'evidenza digitale non nasce tuttavia con l'introduzione dell'\gls{ai}. Le tecniche anti-forensi comprendono, in senso più ampio, strategie finalizzate a ostacolare l'identificazione, l'acquisizione, l'esame o l'interpretazione di dati potenzialmente rilevanti, ad esempio mediante cancellazione selettiva delle tracce, alterazione dei metadati, cifratura, occultamento di file o information hiding. Tra queste tecniche, la steganografia consente di celare informazioni all'interno di contenuti apparentemente ordinari, quali immagini digitali, rendendone più difficile l'individuazione sia da parte dell'analista umano sia attraverso strumenti automatici non specificamente progettati per rilevarla \cite{yaacoub2021digital,sanna2025cybercrime}. Nel presente lavoro, la steganografia viene richiamata come esempio rappresentativo del dominio anti-forense, mentre la valutazione sperimentale si concentra su manipolazioni dell'immagine direttamente pertinenti alla robustezza della classificazione visuale automatizzata.

Parallelamente, la letteratura sull'adversarial machine learning ha dimostrato che i modelli di \gls{ml} e \gls{dl} possono essere vulnerabili a perturbazioni progettate specificamente per indurre errori di classificazione, anche quando tali alterazioni risultano limitate o poco percepibili per un osservatore umano. Questi input manipolati, comunemente indicati come \emph{adversarial examples}, mettono in discussione l'idea che un classificatore ad alte prestazioni su dati puliti sia necessariamente affidabile in scenari reali. La definizione del \emph{threat model}, la distinzione tra livelli differenti di conoscenza dell'attaccante e la necessità di valutazioni orientate alla sicurezza costituiscono pertanto elementi centrali nell'analisi di sistemi automatici esposti a input ostili \cite{goodfellow2015explaining,carlini2017towards,biggio2018wild,nowroozi2021survey}.

Nel dominio forense, attacchi adversariali e tecniche anti-forensi devono essere mantenuti concettualmente distinti. I primi mirano primariamente a sfruttare le fragilità della funzione decisionale di un modello; le seconde possono operare sull'evidenza digitale anche in assenza di qualunque sistema di \gls{ai}, con l'obiettivo più generale di nascondere, degradare o rendere meno informativo un artefatto. Le due prospettive possono tuttavia intersecarsi quando una manipolazione intenzionale riduce la capacità di un sistema AI-based di segnalare contenuti investigativamente rilevanti. In tale scenario, un attacco alla classificazione automatizzata assume anche una finalità anti-forense, poiché può incidere sul triage, sulla prioritizzazione e sulla successiva revisione umana delle evidenze \cite{sanna2024preliminary,sanna2025pixelperfect}.

La nozione di robustezza adottata nella presente tesi deve quindi essere intesa come capacità del sistema di mantenere un comportamento sufficientemente affidabile non soltanto in presenza di dati non manipolati, ma anche rispetto a perturbazioni intenzionali, trasformazioni realistiche dell'immagine e contenuti non pienamente conformi alle condizioni attese di analisi. In questa prospettiva, un sistema non robusto può ridurre la probabilità che immagini rilevanti vengano correttamente portate all'attenzione dell'analista oppure può aumentare il rumore analitico generato da classificazioni errate. La robustezza non costituisce pertanto un problema marginale di ottimizzazione del modello, ma una componente essenziale dell'affidabilità operativa dell'automazione in un workflow forense.

% ─────────────────────────────────────────────────────────────────────────────
\section{Problema di ricerca e rilevanza in scenari forensi realistici}
% ─────────────────────────────────────────────────────────────────────────────
Nonostante l'ampia letteratura sugli attacchi adversariali e sulla robustezza dei modelli di \gls{ml}, rimane ancora limitata la valutazione di tali criticità in scenari forensi realistici. Molti studi si concentrano su benchmark standardizzati, dataset costruiti per finalità generiche di \textit{computer vision} o condizioni sperimentali lontane dai workflow investigativi effettivi. Per contro, i lavori maggiormente aderenti alla pratica forense mostrano che l'analisi automatizzata di immagini deve confrontarsi con grandi raccolte estratte da evidenze digitali, classi operative eterogenee e contenuti rilevanti per specifiche esigenze investigative. In questo senso, Silva et al. propongono un sistema AI-based per classificare immagini contenute in evidenze criminali, includendo categorie quali armi da fuoco, munizioni, documenti, screenshot e nudità; più recentemente, Sanna et al. hanno evidenziato che funzionalità automatiche integrate in software impiegati per l'analisi forense possono mostrare fragilità rispetto a contenuti manipolati o non conformi alle condizioni ordinarie di utilizzo \cite{silva2021microservices,sanna2024preliminary,sanna2025pixelperfect}.

Nel dominio della \gls{df}, il problema non consiste soltanto nello stabilire se un classificatore possa essere indotto in errore, ma nel comprendere in quale misura tale fragilità incida sulla capacità del sistema di supportare decisioni affidabili durante l'analisi dell'evidenza. Ciò richiede set di valutazione aderenti allo scenario operativo, composti da contenuti eterogenei, immagini soggette a trasformazioni plausibili, casi borderline e situazioni in cui la distinzione tra contenuto rilevante, irrilevante o ambiguo non sia sempre netta. Tale esigenza è coerente con la ricerca DFRWS sulla valutazione di sistemi automatici per contenuti visivi: sebbene riferito al \textit{perceptual hashing} e non alla classificazione semantica di armi, il framework PHASER evidenzia la necessità di studiare gli strumenti su dataset costruiti in funzione dello specifico scenario di impiego e dei relativi casi anomali \cite{mckeown2024phaser}. Nella presente tesi, questo principio viene applicato alla valutazione di modelli e software AI-based su contenuti \texttt{weapon}, \texttt{non\_weapon} e casi fuori distribuzione, includendo manipolazioni adversariali e trasformazioni anti-forensi.

In questo quadro, la vulnerabilità del sistema non ha soltanto una rilevanza teorica, ma produce conseguenze operative immediate. Errori di classificazione indotti da input manipolati possono alterare la priorità attribuita alle evidenze, generare rumore analitico o ridurre la capacità del sistema di segnalare contenuti potenzialmente rilevanti. La robustezza, pertanto, non riguarda soltanto la stabilità del modello rispetto a perturbazioni controllate, ma la qualità complessiva del supporto che esso è in grado di offrire all'interno del workflow forense \cite{sanna2024preliminary,sanna2025pixelperfect}.

Il problema di ricerca affrontato in questa tesi si colloca precisamente in tale intersezione tra robustezza algoritmica e affidabilità operativa. L'obiettivo non è limitarsi a osservare una generica degradazione prestazionale, ma analizzare in modo sistematico come differenti modelli di classificazione automatica e differenti strumenti basati su \gls{ai} si comportino quando vengono impiegati in uno scenario forense realistico e sottoposti a perturbazioni adversariali e anti-forensi. Questa prospettiva consente di superare una valutazione puramente astratta della sicurezza dei modelli, orientando l'analisi verso le conseguenze concrete che tali vulnerabilità possono produrre sul triage, sulla selezione delle evidenze e, più in generale, sull'intero workflow investigativo.

La rilevanza della ricerca è dunque duplice. Da un lato, essa è scientifica, poiché contribuisce alla definizione di un quadro sperimentale più aderente alle esigenze del dominio forense e alla comprensione dei limiti dei modelli basati su \gls{ai} in presenza di input ostili. Dall'altro, essa è operativa, poiché fornisce elementi utili per valutare criticamente l'uso di sistemi automatici nei contesti investigativi, promuovendo un approccio che non privilegi la sola accuratezza nominale, ma anche la robustezza, la tracciabilità e la verificabilità del comportamento del modello. In questo senso, la tesi si inserisce in una linea di ricerca che mira a contribuire alla definizione di benchmark più specifici e di protocolli comparativi maggiormente condivisi per la valutazione dell'\gls{ai} in \gls{df} \cite{silva2021microservices,mckeown2024phaser,sanna2024preliminary,sanna2025pixelperfect}.

% ─────────────────────────────────────────────────────────────────────────────
\section{Obiettivi della ricerca}
% ─────────────────────────────────────────────────────────────────────────────
Alla luce delle criticità evidenziate nei paragrafi precedenti, questa tesi si propone di valutare in modo sistematico la robustezza operativa di modelli proxy e software AI-based impiegati, nel contesto della \gls{df}, per la classificazione automatica di immagini di potenziale interesse investigativo. L'obiettivo generale è verificare in quale misura tali sistemi mantengano prestazioni affidabili rispetto alle condizioni di riferimento rappresentate da input non manipolati, quando vengono esposti a immagini degradate, alterate o intenzionalmente perturbate. Particolare attenzione è dedicata sia a manipolazioni riconducibili a strategie adversariali o anti-forensi, sia al comportamento dei sistemi in presenza di contenuti \gls{ood}, potenzialmente in grado di produrre errori operativamente rilevanti. In uno scenario investigativo, tali manipolazioni possono inoltre essere applicate intenzionalmente da soggetti interessati a ridurre la probabilità che contenuti rilevanti vengano individuati mediante procedure automatiche di rilevamento o triage. Recenti studi hanno altresì mostrato che strumenti di generative \gls{ai} general-purpose possono agevolare la produzione di codice destinato all'occultamento di informazioni in immagini digitali mediante tecniche steganografiche; pur non coincidendo con le trasformazioni valutate sperimentalmente nella presente tesi, tale evidenza amplia il quadro delle possibili condotte anti-forensi abilitate o facilitate da strumenti automatici \cite{sanna2024preliminary,sanna2025pixelperfect,sanna2025cybercrime}.

In questa prospettiva, il lavoro si articola attorno a quattro obiettivi principali. Il primo consiste nella definizione di un protocollo sperimentale rigoroso, documentato e riproducibile per l'analisi della robustezza di sistemi di classificazione automatizzata di immagini in un contesto forense realistico. Tale protocollo comprende la costruzione di un dataset finale validato manualmente, la definizione di un sottoinsieme destinato alla valutazione adversariale, la generazione controllata delle perturbazioni e la predisposizione di un bundle di valutazione utilizzabile anche con strumenti forensi commerciali.

Il secondo obiettivo consiste nel confrontare modelli eterogenei, differenti per architettura e paradigma di rappresentazione, al fine di osservare eventuali differenze di comportamento rispetto a immagini clean, contenuti \gls{ood}, manipolazioni anti-forensi e attacchi adversariali. In particolare, i modelli proxy consentono di effettuare un'analisi quantitativa controllata della degradazione prestazionale, mentre gli strumenti forensi AI-based permettono di esaminare la trasferibilità operativa del problema in ambienti applicativi utilizzati nell'analisi digitale.

Il terzo obiettivo riguarda la misurazione dell'impatto delle manipolazioni sulla capacità dei sistemi di riconoscere correttamente immagini appartenenti alle classi di interesse forense. L'analisi non si limita alla variazione dell'accuratezza complessiva, ma considera anche le modalità di errore introdotte, il comportamento rispetto a contenuti fuori distribuzione, la confidenza associata alle decisioni errate e, in chiave diagnostica, l'interpretazione qualitativa delle predizioni mediante tecniche di \gls{xai}.

Il quarto obiettivo consiste nel ricondurre le evidenze sperimentali a una valutazione dell'affidabilità operativa dell'automazione nei workflow di \gls{df}. In tale prospettiva, i sistemi basati su \gls{ai} non vengono considerati strumenti sostitutivi del giudizio dell'analista, bensì supporti al triage e all'esame dei contenuti digitali, la cui affidabilità deve essere verificata in condizioni realistiche e potenzialmente ostili. L'analisi mira pertanto a evidenziare limiti, rischi operativi e requisiti metodologici necessari affinché l'impiego di tali strumenti rimanga compatibile con un approccio human-in-the-loop e con le esigenze di verificabilità proprie dell'attività forense \cite{costantini2019digital,faqir2023digital,sanna2024preliminary}.

L'impostazione adottata non mira quindi soltanto a misurare l'accuratezza di classificatori in un ambiente controllato, ma a comprendere come l'interazione tra modelli, strumenti, dati e perturbazioni possa incidere concretamente sul supporto fornito all'analista. In tal senso, l'obiettivo della ricerca non è esclusivamente prestazionale, ma anche metodologico e applicativo: definire un protocollo di valutazione che consenta di giudicare la robustezza di sistemi basati su \gls{ai} in relazione al loro possibile utilizzo in scenari investigativi realistici, mantenendo centrale la supervisione umana nella lettura e nella validazione dei risultati.


% ─────────────────────────────────────────────────────────────────────────────
\section{Contributo della tesi}
% ─────────────────────────────────────────────────────────────────────────────
Il contributo principale della presente tesi risiede nella definizione e nell'applicazione di una metodologia sperimentale per la valutazione della robustezza operativa di sistemi AI-based destinati alla classificazione automatizzata di immagini in un contesto di \gls{df}. A differenza di studi focalizzati prevalentemente su benchmark standard di \textit{computer vision} o su valutazioni di robustezza condotte in domini generici, il presente lavoro adotta una prospettiva orientata alle esigenze dell'analisi forense: non considera soltanto la capacità nominale di classificazione, ma esamina anche la tenuta del sistema in presenza di input ostili, degradati, ambigui o fuori distribuzione \cite{nowroozi2020phd,hendrycks2019benchmarking,sanna2024preliminary,sanna2025pixelperfect}.

Il protocollo sperimentale adottato nella tesi è denominato \gls{fairlab}. Esso organizza in modo documentato e riproducibile le componenti necessarie alla valutazione: dataset finale congelato, modelli proxy, perturbazioni adversariali, trasformazioni anti-forensi, software commerciali selezionati per la valutazione operativa, metriche quantitative, analisi di \gls{xai} e controlli di tracciabilità. \gls{fairlab} non viene proposto come nuovo strumento software autonomo né come architettura teorica generale, ma come protocollo metodologico utilizzato per rendere trasparente, verificabile e riproducibile la valutazione sperimentale svolta nella presente tesi.

Un primo contributo consiste nell'integrazione, all'interno del medesimo disegno sperimentale, di due livelli di valutazione complementari. Il primo riguarda la valutazione controllata dei modelli proxy, attraverso la quale dataset, perturbazioni, metriche e analisi di \gls{xai} possono essere esaminati in condizioni riproducibili e comparabili. Il secondo riguarda la valutazione operativa black-box dei software commerciali selezionati, condotta mediante un insieme comune di immagini organizzato per condizioni sperimentali e utilizzato per osservare il comportamento degli output disponibili all'analista. Tale integrazione consente di collegare l'analisi quantitativa dei modelli alla verifica del comportamento di software effettivamente utilizzabili in attività di analisi di contenuti digitali, senza assumere accesso alla logica interna dei relativi componenti AI-based \cite{costantini2019digital,sanna2024preliminary,sanna2025pixelperfect}.

Un secondo contributo riguarda la costruzione di un dataset finale congelato, nonché la definizione di un protocollo di selezione coerente con il problema investigativo affrontato. Il dataset finale è organizzato in tre classi semanticamente distinte, denominate \texttt{weapon}, \texttt{non\_weapon} e \texttt{ood}, per un totale di 1500 immagini bilanciate. La classe \texttt{weapon} include esclusivamente armi individuali realistiche e chiaramente visibili; la classe \texttt{non\_weapon} comprende negativi realistici coerenti con il task; la classe \texttt{ood} raccoglie casi borderline o fuori distribuzione rispetto al task principale, quali coltelli, spade, repliche, armi giocattolo o airsoft, render sintetici, immagini tratte da videogiochi, veicoli o armamenti non individuali, scene di guerra, esplosioni, immagini degradate o contenuti anomali rispetto al dominio di interesse. A partire dal dataset finale è inoltre definito un subset binario di 1000 immagini, composto dalle sole classi \texttt{weapon} e \texttt{non\_weapon}, destinato alla generazione e alla valutazione delle perturbazioni adversariali e delle trasformazioni anti-forensi. La selezione dei campioni segue un paradigma \textit{human-in-the-loop}, nel quale la componente manuale non viene occultata, ma documentata come parte integrante del protocollo metodologico, rendendo espliciti criteri decisionali, assunzioni e possibili limiti della costruzione del corpus sperimentale \cite{sambasivan2021data}.

Un terzo contributo riguarda la definizione di un piano di valutazione fondato su dati eterogenei, criteri di classificazione coerenti con il problema investigativo affrontato e un insieme di manipolazioni selezionate per rappresentare due prospettive distinte ma operativamente correlate: perturbazioni adversariali progettate per alterare la risposta del classificatore e trasformazioni anti-forensi plausibili sul piano del ciclo di vita e dell'occultamento dell'immagine digitale. In questo modo, la tesi non si limita a verificare se un sistema possa essere indotto in errore, ma analizza anche la natura dell'errore e le sue possibili implicazioni sul triage e sulla selezione preliminare di contenuti rilevanti \cite{barni2018cf,nowroozi2024verifying,yaacoub2021digital}.

Un quarto contributo consiste nella comparazione tra modelli differenti per architettura e paradigma di rappresentazione, unitamente alla valutazione black-box dei software commerciali inclusi nel protocollo. Tale impostazione consente di evitare conclusioni esclusivamente dipendenti da un singolo classificatore e di osservare se le fragilità riscontrate in condizioni controllate assumano rilevanza anche rispetto agli output disponibili in software utilizzati nel contesto dell'analisi forense. L'analisi di \gls{xai}, applicata ai modelli proxy, svolge in questo quadro una funzione diagnostica complementare: non certifica la correttezza della classificazione, ma supporta l'interpretazione qualitativa di casi rappresentativi di errore o instabilità.

Nel complesso, la tesi propone una lettura della robustezza non soltanto come proprietà tecnica del modello, ma come requisito funzionale per l'impiego controllato dell'automazione in \gls{df}. L'analisi empirica è pertanto finalizzata a valutare in quale misura dati fuori distribuzione, perturbazioni adversariali e trasformazioni anti-forensi incidano sulla stabilità delle classificazioni e sulla qualità del supporto offerto all'analista, mantenendo centrale la supervisione umana nell'interpretazione e nella verifica dei risultati.


% ─────────────────────────────────────────────────────────────────────────────
\section{Rilevanza tecnico-forense e profili di affidabilità probatoria}
% ─────────────────────────────────────────────────────────────────────────────
In ambito forense, la correttezza tecnica di uno strumento non può essere valutata indipendentemente dal contesto in cui esso viene utilizzato e dalle conseguenze che i suoi risultati possono produrre sull'analisi dell'evidenza. L'introduzione di sistemi basati su \gls{ai} nei workflow investigativi solleva pertanto una questione che è al tempo stesso tecnica e forense: fino a che punto un output automatizzato può essere considerato sufficientemente affidabile da orientare il lavoro dell'analista, senza compromettere la completezza dell'esame o introdurre errori sistematici difficilmente rilevabili \cite{faqir2023digital,sanna2025cybercrime}.

Il problema assume rilievo particolare quando i sistemi automatici operano in modalità sostanzialmente opaca, ovvero quando l'analista dispone del risultato della classificazione ma non di una piena comprensione del processo decisionale che lo ha generato. In tali casi, la fragilità del modello rispetto a input manipolati o inattesi non rappresenta soltanto un limite algoritmico, ma può riflettersi sulla qualità dell'intero processo investigativo. Un falso negativo può determinare la mancata selezione di un contenuto rilevante; un falso positivo può generare rumore, deviare risorse e appesantire la verifica manuale. In questo quadro, gli approcci di \gls{xai} possono fornire un supporto diagnostico all'analista, consentendo di esaminare, con le necessarie cautele, quali componenti dell'input abbiano maggiormente contribuito alla predizione e se il comportamento del modello risulti coerente o anomalo in presenza di immagini manipolate. Tali approcci non rendono automaticamente affidabile la classificazione né sostituiscono la validazione sperimentale o il controllo umano, ma contribuiscono a documentare e interpretare criticamente il ruolo svolto dall'automazione nell'analisi forense \cite{Adadi2018,Doshi-Velez2017}.

Da questo punto di vista, la robustezza operativa si collega direttamente all'esigenza di tracciabilità, verificabilità e controllo critico che caratterizza la \gls{df}. Un sistema automatico può certamente costituire un supporto efficace per il triage e la prioritizzazione delle evidenze, ma il suo impiego richiede che i risultati siano interpretati alla luce dei limiti del modello, delle condizioni sperimentali in cui è stato valutato e dei possibili margini di errore. La presente tesi si colloca precisamente in questo spazio problematico: non assume l'automazione come intrinsecamente affidabile, ma ne esamina la tenuta empirica per comprendere in quali condizioni essa possa offrire un supporto utile e in quali, invece, richieda particolare cautela.

Sotto il profilo probatorio, tale impostazione non implica che il modello debba essere considerato fonte autonoma di verità tecnica, bensì che il suo eventuale impiego debba essere accompagnato da un'adeguata consapevolezza metodologica. In altri termini, la rilevanza del presente studio risiede anche nel fornire elementi utili a distinguere tra uso meramente strumentale dell'\gls{ai} e affidamento acritico su decisioni algoritmiche non sufficientemente robuste. Ciò risulta particolarmente importante in un contesto in cui la prova digitale richiede procedure rigorose e in cui l'adozione di strumenti automatici deve poter essere difesa non solo sul piano dell'efficienza, ma anche su quello della controllabilità tecnica e della loro affidabilità in ambito forense \cite{cpp1988,aiact2024,swgde_guidelines}.

% ─────────────────────────────────────────────────────────────────────────────
\section{Struttura della tesi}
% ─────────────────────────────────────────────────────────────────────────────
La tesi è articolata in sei capitoli, secondo una progressione che dal contesto teorico e dallo stato dell'arte conduce alla metodologia sperimentale, all'analisi dei risultati e alla discussione delle relative implicazioni operative.

Il Capitolo~\ref{chap:introduction} introduce il problema di ricerca, gli obiettivi del lavoro e il contributo della tesi rispetto alla robustezza operativa di sistemi AI-based impiegati per la classificazione automatizzata di immagini in \gls{df}.

Il Capitolo~\ref{chap:background} presenta i fondamenti teorici relativi alla prova digitale, all'impiego dell'\gls{ai} nell'analisi forense, alla classificazione automatica di immagini, all'adversarial machine learning, alle tecniche anti-forensi e alla \gls{xai}.

Il Capitolo~\ref{chap:state-of-the-art} analizza criticamente la letteratura rilevante, individuando il gap relativo alla valutazione integrata della robustezza di modelli e software AI-based in scenari forensi realistici.

Il Capitolo~\ref{chap:methodology} descrive il protocollo sperimentale FAIR-Lab, includendo la costruzione e validazione del dataset, i modelli proxy, le perturbazioni considerate, la valutazione black-box dei software selezionati, le metriche e i controlli di riproducibilità.

Il Capitolo~\ref{chap:results} presenta i risultati sperimentali, discutendo il comportamento dei sistemi su immagini non manipolate, contenuti fuori distribuzione, trasformazioni anti-forensi, perturbazioni adversariali e casi qualitativi esaminati mediante \gls{xai}.

Il Capitolo~\ref{chap:conclusions} sintetizza le evidenze principali, ne discute le implicazioni per l'impiego controllato dell'automazione in ambito forense e individua i limiti dello studio e le possibili direzioni di ricerca futura.
